\subsubsection{Unique witnesses with almost the entire gap as separation}
\label{sec:unique-orbits}

Larger prime-order orbits strengthen the separation while preserving
uniqueness and efficient witness recovery. The evaluation domain is a
union of scaled small orbits, with a bounded number of extra coordinates.
It is not a prescribed multiplicative subgroup.

\begin{lemma}[Rigidity of root sums on scaled prime-order orbits]
\label{ou:root-sums}
Let $d$ be prime, $p\equiv1\pmod d$, and let $\omega\in\F_p$ have
order $d$. Fix $m\ge2$ and $0\le c<d$, and put
\[
 A=2^{m+1}(d+2^c).
\]
If $p>A^{d-1}$, the following points are distinct:
\[
 \{2^i\omega^j:1\le i\le m,\ 0\le j<d\}
 \ \cup\ \{2^{m+t}:1\le t\le c\}
 \ \cup\ \{\omega^j:0\le j<d\}.
\]
Every zero-sum subset of the first two sets is a union of complete
$d$-element orbits from the first set; it contains no extra coordinate.
\end{lemma}
\begin{proof}
Let $\zeta$ be a complex primitive $d$th root of unity. The polynomial
$\Phi_d(X)=1+X+\cdots+X^{d-1}$ is irreducible over $\mathbb Q$ by Eisenstein
after translating $X$ by one. For an integer polynomial $P$ of degree
at most $d-1$, if $P(\omega)=0$ but $P(\zeta)\ne0$, its nonzero
integer resultant with $\Phi_d$ is divisible by $p$. On the other hand,
if the sum of the absolute values of its coefficients is at most $A$,
\[
 |\operatorname{Res}(\Phi_d,P)|
 =\prod_{j=1}^{d-1}|P(\zeta^j)|\le A^{d-1}<p,
\]
a contradiction.

A difference between two distinct displayed formal points is nonzero
over $\mathbb C$ and has coefficient sum less than $A$, so their reductions
cannot collide. Now lift a zero-sum subset of nonpadding points to
$P(\zeta)$. Its coefficient sum is less than $A$, hence $P(\zeta)=0$.
Its degree is at most $d-1$, so all its coefficients must be equal.
Each coefficient records a sum of distinct powers of two. Uniqueness
of binary expansion forces the same selected indices in every orbit
position. Thus each core orbit is either fully selected or absent,
and an extra coordinate, which occurs only in position zero, is absent.
\end{proof}

\begin{theorem}[Explicit orbit lines with unique witnesses]
\label{ou:finite}
With the preceding notation, take $1\le D<m$ and assume also
\[
 p>2^{dD(2m-D+1)/2}.
\]
There is a length-$n=d(m+1)+c$, dimension-$K=dD-1$ Reed--Solomon
code and a received line with exactly $\binom mD$ nearby parameters
at radius
\[
 \theta=1-\frac Kn-\frac{d+1}{n}.
\]
Every nearby point has a unique codeword, recoverable in polynomial
time in $\log p$. The exceptional point has exact distance
$1-K/n-1/n$; every nearby point has exact distance $\theta$.
The direction has weight $d$, attaining the distance-improvement bound.
There is no correlated agreement at the tested threshold.
Given $\omega$, construction and decoding are deterministic. An order-$d$
root can be found by a zero-error randomized procedure in expected
polynomial time in $\log p$.
\end{theorem}
\begin{proof}
Use the three displayed sets as the core, extra coordinates, and padding,
respectively. For $D$-subsets $I\subseteq\{1,\ldots,m\}$ put
\[
 H_I(X)=\prod_{i\in I}(X^d-2^{di}).
\]
Fix a reference $w=H_{I_0}$, let $f$ be its evaluation word, and let $g$
equal one on the padding and zero elsewhere. The polynomial $w-H_I$
has degree at most $dD-d<K$ and supplies the nearby parameter
\[
 z_I=-H_I(1)=(-1)^{D+1}\prod_{i\in I}(2^{di}-1).
\]
It matches $dD=K+1$ core coordinates and all $d$ padding coordinates.
The root bound gives exact far agreement $K+1$ for $f$; changing only
$d$ coordinates proves the exact nearby distance.

Conversely, any nearby codeword must match all padding coordinates
and $K+1=dD$ nonpadding coordinates. Its difference from $w$ is monic
of degree $dD$, has zero coefficient of $X^{dD-1}$, and is the locator
of those $dD$ distinct roots. Their sum is zero, so
Lemma~\ref{ou:root-sums} forces them to be exactly $D$ complete core
orbits. Thus every nearby codeword belongs to the displayed family.

The positive integer products are less than
$2^{dD(2m-D+1)/2}<p$. The proof and greedy decoder of
Lemma~\ref{pd:integer-products} apply with base $2^d$ in place of four.
The normalized product is still greater than $2/3$, and the product of
all later factors exceeds the membership threshold. This proves
injectivity, the exact count, and unique witness recovery from the
signed integer lift of $z$.
Finally, sampling $h\in\F_p^*$ and testing
$h^{(p-1)/d}\ne1$ finds $\omega$ in expected $d/(d-1)$ trials.
Every accepted root has order $d$, so this step has no correctness
failure event.
\end{proof}

\begin{corollary}[Unambiguous numerical failure with nearly full separation]
\label{ou:asymptotic}
Fix a rational rate $\rho=a/b\in(0,1)$ in lowest terms and a prime
$d\nmid a$. For every sufficiently large prime $p\equiv1\pmod d$,
the construction has exact rate $\rho$, length
$n=\Theta_{\rho,d}(\sqrt{\log p})$, and gap $\eta=(d+1)/n$, with
\[
 \dist(f,C)=\theta+\frac d{d+1}\eta.
\]
Every nearby point has a unique, efficiently recoverable witness.
The radius is strictly below Elias and, for every fixed $c_1>0$ and
$0<c_2<1+1/d$,
\[
 \frac{\binom mD}{c_1n2^{c_2H_2(\rho)/\eta}}
 =2^{\Omega_{\rho,d,c_2}(n)}.
\]
In particular, for every fixed $\varepsilon>0$, infinitely many prime
fields admit such counterexamples to the $c_2=1$ prescription, with separation at least
$(1-\varepsilon)\eta$ and no decoding ambiguity anywhere on the nearby set.
\end{corollary}
\begin{proof}
Take $n=bt$ with $at\equiv-1\pmod d$, let $c$ be the residue of $n$
in $\{0,\ldots,d-1\}$, and put
\[
 m=(n-c)/d-1,\qquad D=(\rho n+1)/d.
\]
Choosing $n$ near a sufficiently small constant multiple of
$\sqrt{\log_2p}$ ensures both field bounds, with $1\le D<m$ eventually.
Then
\[
 \log_2\binom mD=\frac{n}{d}H_2(\rho)+O_{\rho,d}(\log n).
\]
Comparing this with $c_2nH_2(\rho)/(d+1)$ proves the ratio. Since
$\eta\log_2p\to\infty$, strict Elias holds. Choosing a sufficiently
large prime $d$ not dividing $a$ gives the final separation claim.
There are infinitely many primes congruent to one modulo this $d$.
The gap still tends to zero; no fixed-positive-gap conclusion is asserted.
\end{proof}

Exact half-rate certificates illustrate the tradeoff:
\begin{center}
\begin{tabular}{@{}rrrrll@{}}
\toprule
Prime & $d$ & $n$ & $K$ & Separation & Ratio $>$\\
\midrule
$2^{9689}-1$&3&274&137&$3\eta/4$&$2^9$\\
$2^{19937}-1$&5&518&259&$5\eta/6$&$2^2$\\
$2^{44497}-1$&7&908&454&$7\eta/8$&$1$\\
\bottomrule
\end{tabular}
\end{center}
Here the ratio compares the exact nearby count $\binom mD$ with
$n2^{1/\eta}$. The respective $(m,D,c)$ are $(90,46,1)$,
$(102,52,3)$, and $(128,65,5)$. The accompanying integer certificates
verify the field bounds, primality, strict Elias, and these comparisons.
Unlike the complete-coverage examples, these lines have a sparse nearby
bank. Their witness classification and decoder have no sampling failure.

\paragraph{Smaller finite fields using roots of two.}
When the field also contains a $d$th root of two, the same decoder can
use the factors $2^i-1$ instead of $2^{di}-1$. This reduces the integer
product bound substantially. The following statement requires both
specified roots in the field; it does not replace the prime-progression
hypothesis of Corollary~\ref{ou:asymptotic}.

\begin{proposition}[A root-of-two refinement]\label{ou:kummer}
Let $d$ be prime and let $\omega,\alpha\in\F_p$ satisfy
$\operatorname{ord}(\omega)=d$ and $\alpha^d=2$.
Take $m\ge2$, $1\le D<m$, $0\le c<d$, and put
\[
 n=d(m+1)+c,\qquad K=dD-1,\qquad
 A=n2^{\lceil(m+c+1)/d\rceil}.
\]
Assume
\[
 p>A^{d(d-1)},\qquad p>2^{D(2m-D+3)/2}.
\]
Use the core $\{\alpha^i\omega^j:2\le i\le m+1,\ 0\le j<d\}$,
extra coordinates $\alpha^{m+2},\ldots,\alpha^{m+c+1}$, and padding
$\{\omega^j:0\le j<d\}$.
These points are distinct. The corresponding locator line has exactly
$\binom mD$ nearby parameters at radius $1-K/n-(d+1)/n$;
each has a unique codeword, recoverable in polynomial time in $\log p$.
Parameter zero has exact distance $1-K/n-1/n$.
The construction and decoder are deterministic given the two roots.
\end{proposition}
\begin{proof}
Let $u=2^{1/d}>0$ and let $\zeta$ be a complex primitive $d$th root.
The field $L=\mathbb Q(u,\zeta)$ has degree $d(d-1)$: the degrees
$d$ and $d-1$ of its two subfields are relatively prime, so their
product both divides and bounds $[L:\mathbb Q]$.
The ring $\mathbb Z[u,\zeta]$ is free on the basis
$u^s\zeta^j$ for $0\le s<d$, $0\le j<d-1$.
The given finite-field roots define a ring homomorphism to $\F_p$.
Consequently a nonzero element mapping to zero has nonzero integer
norm divisible by $p$, as follows by reducing its multiplication
matrix modulo $p$.

For a lifted subset sum or a difference of two displayed formal
coordinates, every conjugate has absolute value at most $A$.
Thus its norm has absolute value at most $A^{d(d-1)}<p$.
A modular zero must therefore already be zero in $L$. Distinct formal
coordinates have different positive magnitudes unless their indices
agree, in which case their root-of-unity factors are distinct. This
proves distinctness after reduction.

For a subset sum of nonpadding coordinates, group exponents of $u$
according to their residue modulo $d$, using $u^{qd+s}=2^qu^s$.
The powers $1,u,\ldots,u^{d-1}$ are linearly independent over
$\mathbb Q(\zeta)$, so each resulting polynomial in $\zeta$ vanishes.
Its degree is at most $d-1$, forcing all its coefficients to agree.
Within each residue class, each coefficient is a sum of distinct
powers of two. Binary uniqueness again forces whole core orbits and
excludes all extra coordinates.

The all-codeword classification from Theorem~\ref{ou:finite} now
applies, with locators
\[
 H_I(X)=\prod_{i\in I}(X^d-2^i),\qquad
 I\subseteq\{2,\ldots,m+1\},\quad |I|=D.
\]
Their signed labels lift to positive integers
$T_I=\prod_{i\in I}(2^i-1)<2^{D(2m-D+3)/2}<p$.
Writing $s=\sum_{i\in I}i$, the normalized product satisfies
\[
 \tfrac12<\prod_{i\in I}(1-2^{-i})<1,
\]
so $T_I$ has exactly $s$ binary digits. At index $i$, membership is
equivalent to the remaining normalized product being at most
$1-2^{-i}$; the product of all later factors is strictly greater than
that threshold. The greedy decoder therefore works in base two.
Starting the indices at two is essential to this normalization argument.
The remaining degree, count, and distance conclusions are unchanged.
\end{proof}

For a Mersenne prime $p=2^b-1$ with $d\ne b$, one can take
$\alpha=2^v$, where $dv\equiv1\pmod b$. The check
$\alpha^d=2$ is immediate. An order-$d$ root is still required, so
$d\mid p-1$ remains a separate condition.

Exact certificates give the following specific domains. Rates are $K/n$
and need not be one half. Each comparison uses that row's exact binary
entropy, evaluated through integer powers rather than floating-point logs.
\begin{center}
\begin{tabular}{@{}rrrrll@{}}
\toprule
Prime & $d$ & $n$ & $K$ & Separation & Ratio $>$\\
\midrule
$2^{521}-1$&2&80&29&$2\eta/3$&$2^3$\\
$2^{1279}-1$&3&198&68&$3\eta/4$&$2^4$\\
$2^{9689}-1$&5&905&324&$5\eta/6$&$2^{14}$\\
$2^{23209}-1$&13&3588&1338&$13\eta/14$&$2^2$\\
$2^{44497}-1$&19&7068&2830&$19\eta/20$&$1$\\
\bottomrule
\end{tabular}
\end{center}
The ratio is $\binom mD/(n2^{H_2(K/n)/\eta})$, with
$\eta=(d+1)/n$. The respective $(m,D,c)$ are
$(39,15,0)$, $(65,23,0)$, $(180,65,0)$, $(275,103,0)$,
and $(371,149,0)$.
All five primes have independent Lucas--Lehmer replays; saved roots,
domain distinctness, the two field bounds, strict Elias, and representative
witness evaluations are checked independently. These are deterministic
specific-domain certificates once the certified roots are supplied.
The final row has a uniquely recoverable nearby witness at every accepted
parameter, even though the far point is $19/20$ of the capacity gap
outside the radius and the numerical prescription is too small.

\paragraph{Shorter domains by randomizing the orbit representatives.}
The preceding decoder uses the integer size of the binary-scaled
coordinates. If efficient recovery at an arbitrary supplied parameter
is not required, random representatives give shorter domains and a
larger nearby bank relative to the field.

\begin{theorem}[Unique orbit lines of logarithmic length]
\label{ou:random}
Fix a rational rate $\rho=a/b\in(0,1)$ in lowest terms, a prime
$d\nmid a$, and $0<\alpha<1$. For every sufficiently large prime
$p\equiv1\pmod d$, there are exact-rate Reed--Solomon codes of length
\[
 n=\alpha\log_2p+O_{\rho,d}(1)
\]
and received lines with exactly
\[
 J=2^{(H_2(\rho)/d+o(1))n}
\]
nearby parameters at radius $\theta=1-\rho-(d+1)/n$.
Every nearby point has a unique codeword and exact distance $\theta$;
parameter zero has exact distance $1-\rho-1/n$.
Consequently its separation is $d/(d+1)$ of the capacity gap.
The radius is strictly below Elias, correlated agreement fails, and
for every fixed $c_1>0$ and $0<c_2<1+1/d$,
\[
 \frac{J}{c_1n2^{c_2H_2(\rho)/((d+1)/n)}}=p^{\Omega_{\rho,d,\alpha,c_2}(1)}.
\]
A randomized procedure constructs the code and line in expected
polynomial time in $\log p$, with failure probability at most
$p^{\alpha-1+o(1)}$. This theorem does not assert efficient witness
recovery from an arbitrary nearby parameter.
\end{theorem}
\begin{proof}
Choose $n,c,m,D$ by the exact-rate congruences in
Corollary~\ref{ou:asymptotic}, now with $n$ near $\alpha\log_2p$.
Fix an order-$d$ root $\omega$. Sample independent uniform
$a_1,\ldots,a_m,b_1,\ldots,b_c\in\F_p$ and use the domain
\[
 \{a_i\omega^j:1\le i\le m,\ 0\le j<d\}
 \ \cup\ \{b_1,\ldots,b_c\}
 \ \cup\ \{\omega^j:0\le j<d\}.
\]
The probability of a collision among its $n$ formal points is at most
$\binom n2/p$.

For $p>d^{d-1}$, no nonempty proper subset of
$\{1,\omega,\ldots,\omega^{d-1}\}$ sums to zero. Indeed, its lift to
$\mathbb Z[\zeta_d]$ is nonzero by irreducibility of $\Phi_d$, and
its nonzero norm has absolute value at most $d^{d-1}<p$.
Thus every subset of the nonpadding domain other than a union of
complete core orbits has a sum that is a nonzero linear form in the
sampled variables. The probability that any such subset has sum zero
is at most
\[
 \frac{2^{dm+c}}p=\frac{2^{n-d}}p.
\]
For two distinct $D$-subsets $I,J$, the polynomial
\[
 \prod_{i\in I}(1-a_i^d)-\prod_{j\in J}(1-a_j^d)
\]
is nonzero and has degree at most $dD$. The probability of any product
collision is therefore at most
\[
 \frac{dD}{p}\binom{\binom mD}{2}.
\]
On every output avoiding these three bad events, the root-locator
classification in Theorem~\ref{ou:finite} applies verbatim: every
nearby codeword comes from exactly $D$ whole orbits, and the distinct
products make all these codewords occur at distinct parameters.

Before conditioning on a distinct domain, the total bad probability
is at most
\begin{equation}\label{eq:ou-random-failure}
 B=\frac{\binom n2+2^{n-d}+dD\binom{\binom mD}{2}}p.
\end{equation}
This is $p^{\alpha-1+o(1)}$, since
$\binom mD^2\le 2^{2m}$ and $2m\le n-d$ for $d\ge2$.
Rejecting invalid domains changes the bound by at most the factor
$(1-\binom n2/p)^{-1}=1+o(1)$.
The same entropy calculation proves the numerical ratio, and
$((d+1)/n)\log_2p\to(d+1)/\alpha$ gives strict Elias with room to spare.
The sampler checks domain validity and forms the reference polynomial;
it does not enumerate subsets or test the two combinatorial good events.
\end{proof}

For example, the exact union bound gives the following half-rate rows.
The ratio compares $J$ with $n2^{1/\eta}$; the probability column bounds
failure after rejection sampling for distinct coordinates.
\begin{center}
\begin{tabular}{@{}rrrrlll@{}}
\toprule
Prime & $d$ & $n$ & $K$ & Separation & Ratio $>$ & Failure $<$\\
\midrule
$2^{521}-1$&3&466&233&$3\eta/4$&$2^{24}$&$2^{-57}$\\
$2^{521}-1$&5&468&234&$5\eta/6$&$2$&$2^{-57}$\\
$2^{1279}-1$&7&1146&573&$7\eta/8$&$2^4$&$2^{-139}$\\
\bottomrule
\end{tabular}
\end{center}
These are guarantees for the sampling procedure. The arithmetic
certificates do not deterministically certify the combinatorial good
events for a particular sampled domain.

